
\subsubsection{Prerequisites}
Prerequisites for a well functioning installation of \edps\ and \edps-gui\ are:

\begin{itemize}
\item    Recent Firefox or Chrome browser, Python 3.11 or higher (but there are issues with Python 3.14).

\item At least one ESO pipeline with EDPS workflow should be in your
  system. To install the desired ESO pipelines, follow the
  instructions in the ESO pipelines pages. NOTE: the {\tt apptainer}
  installation method is currently not supported. After the
  installation, the esorex command must be in the path. To test
  whether the installation was successful, type
\begin{verbatim}
  esorex --recipes 
\end{verbatim}

A list of available recipes should appear.

\item Install {\tt graphviz}, {\tt fv}, and {\tt ds9}, which have to
  be included in the system path (definying aliases not
  enough). On linux, {\tt Graphviz} can be easily installed via:

\begin{verbatim}
 sudo apt install graphviz (Debian, Ubuntu)
 sudo dnf install graphviz (Fedora)
\end{verbatim}
Check the , {\tt Graphviz} webpage for installation instructions for other OS.

{\tt fv} and {\tt ds9}, are optional. To install them, follow the
instructions in corresponding webpages. You can test whether these
three packages are installed and their path are correctly set by
typing on a terminal:

\begin{verbatim}
 dot -V
 fv  -version
 ds9 -version
\end{verbatim}
\end{itemize}

\subsubsection{Installation steps}

To install EDPS follow these steps:

\begin{itemize}

\item Create a new Python virtual environment and activate it:
\begin{verbatim}
  python3 -m venv edpsgui
  . edpsgui/bin/activate
\end{verbatim}
Make sure the python3 version is 3.11 or higher, but not 3.14.

\item Install the required packages:
\begin{verbatim}
 pip install --extra-index-url \
     https://ftp.eso.org/pub/dfs/pipelines/repositories/stable/src \
     edps edpsgui edpsplot adari_core 
\end{verbatim}
\end{itemize}

To run the \edpsgui\ type from a terminal (with the active environment):
\begin{verbatim}
 edps-gui
\end{verbatim}

{\bf Important note}. The first time \edpsgui\ is executed, you will
be asked to specify the directory where the reduction products (fits
files and quality plots) will be stored. The default location is {\tt
  \$HOME/EDPS\_data}. During the first execution, a configuration file
named {\tt application.properties} will also be saved in the directory
(newly created) {\tt \$HOME/.edps}.
